\section{Conclusions and Future Research}

In this paper, we aimed to topologically characterize the grokking phenomenon through comprehensive experiments across two datasets. Our observations indicate that generalization is accompanied by a structural simplification of the model's representation space. Specifically, topological invariants such as homologies and intrinsic dimensions converge to lower values, suggesting the emergence of better-behaved representation manifolds. In these scenarios, generalization can be understood as a process where the network geometrically manipulates its representation manifold to achieve a more simplified, organized state.

Directions for future research include observing this behavior across different cyclic groups $\mathbb{Z}_n$, where $n \in \mathbb{Z}^+$. Additionally, leveraging the natural isomorphism to the n-th roots of unity to observe how it alters the representation space would be highly beneficial. These future experiments could be readily implemented using tokenized datasets or even simplified numerical paired data.